\section*{Bab \ref{ch:g3:food-chains} --- Rantai Makanan}

\begin{solution}{exo:g3:food-chains:1}
``\omterm{def:g3:food-chains:chain}{Dimakan oleh}''. Ia menunjuk dari yang dimakan ke yang memakan ---
mengikuti makanannya.
\end{solution}

\begin{solution}{exo:g3:food-chains:2}
Sebuah \omterm{def:g1:plants-around-us:plant}{tumbuhan} (atau \omterm{def:g1:living-or-not:living}{makhluk hidup} lain yang membuat makanannya
sendiri). Karena hanya \omterm{def:g1:plants-around-us:plant}{tumbuhan} yang makan tanpa memakan siapa pun ---
dari cahaya, air dan udara; setiap mata rantai \omterm{def:g1:animals-around-us:animal}{hewan} hidup dari mata
rantai hijau yang pertama itu, langsung atau melalui yang lain.
\end{solution}

\begin{solution}{exo:g3:food-chains:3}
Rumput $\rightarrow$ belalang $\rightarrow$ kadal $\rightarrow$
elang.
\end{solution}

\begin{solution}{exo:g3:food-chains:4}
Anak panahnya terbalik: rumput $\rightarrow$ kelinci $\rightarrow$
rubah --- rumput \omterm{def:g3:food-chains:chain}{dimakan oleh} kelinci, yang \omterm{def:g3:food-chains:chain}{dimakan oleh} rubah.
\end{solution}

\begin{solution}{exo:g3:food-chains:5}
\omterm{def:g1:plants-around-us:parts}{Daun} pohon ek $\rightarrow$ \omterm{ex:g2:animals-grow-up:butterfly}{ulat} $\rightarrow$ \omterm{prop:g3:sorting-living-things:groups}{burung} gelatik biru
$\rightarrow$ alap-alap.
\end{solution}

\begin{solution}{exo:g3:food-chains:6}
Seladanya \omterm{def:g3:food-chains:chain}{dimakan oleh} siputnya; landaknya memakan siputnya.
\end{solution}

\begin{solution}{exo:g3:food-chains:7}
Misalnya: \omterm{def:g1:plants-around-us:plant}{tumbuhan} air $\rightarrow$ \omterm{ex:g2:animals-grow-up:frog}{berudu} $\rightarrow$ \omterm{prop:g3:sorting-living-things:groups}{ikan} duri
$\rightarrow$ kuntul.
\end{solution}

\begin{solution}{exo:g3:food-chains:8}
Karena makanan dari makanan dari makanannya adalah rumput: belalang
yang lebih sedikit pada rumput yang mencokelat berarti kadal yang
lebih sedikit, dan kadal yang lebih sedikit berarti perburuan yang
kurus bagi elangnya. Sebuah perubahan pada mata rantai pertama
merambat naik sepanjang seluruh rantainya.
\end{solution}

\begin{solution}{exo:g3:food-chains:9}
Meracuni siputnya membuat pemakannya kelaparan: landaknya pergi atau
mati. Kemudian, dengan landaknya --- musuh utama siputnya --- yang
sudah lenyap, siput telanjang dan siput yang selamat berlipat ganda
lebih keras daripada sebelumnya: menyingkirkan mata rantai merambat
baik naik sepanjang rantainya (landak yang lapar) maupun turun
kembali (siput yang tidak diburu).
\end{solution}

\begin{solution}{exo:g3:food-chains:10}
Benar --- melalui bagian tengah rantainya. Kadalnya memakan belalang,
dan belalang memakan rumput: kalau rumputnya gagal, kegagalannya
merambat naik, mata rantai demi mata rantai, dari rumput ke belalang
lalu ke kadal. Kadalnya tidak pernah menyentuh rumput, dan tetap saja
bergantung padanya.
\end{solution}
